\begin{textsample}{NEG Dim 4 – persona\_gpt – Score -1.00 – t115\_gpt.txt}  \label{ex:f4_neg_032}
Across Kenya, smallholder farmers are standing up to laws that criminalise the very practices that have sustained their communities for generations.
Under current Kenyan legislation, sharing, exchanging, or selling uncertified \textbf{seeds} is prohibited and can be punished with up to two years in prison or a KES 1,000,000 fine.
These rules favour multinational corporations and deepen the injustices of a broken global food system, where corporate interests are placed above the rights, knowledge, and resilience of local farmers.
At the heart of this resistance are people who refuse to accept a future where farmers lose control over their \textbf{seeds}.
In Kenya, Nasike is leading Greenpeace Africa ’s Seed is Sovereign campaign, which defends \textbf{seed} sovereignty and backs the court case challenging these unconstitutional \textbf{seed} laws. \textbf{Seed} is Sovereign centres farmers ’ rights to save, share, and exchange \textbf{seeds} as a foundation of food security, biodiversity, and cultural heritage.
By supporting Kenyan smallholder farmers in court and amplifying their voices, the campaign confronts a system that punishes traditional practices while rewarding corporate \textbf{seed} ownership.
This struggle is not isolated to Kenya.
Across the continent, activists are drawing strength and lessons from each other ’s victories.
In Cameroon, Darcise Dolorès, an environmental law student and Greenpeace volunteer, is using the example of the South African rooibos case to help protect Indigenous communities ’ access to palm oil.
That case, in which Indigenous knowledge over rooibos was recognised in South Africa, has become a reference point for how local communities can assert their rights over traditional resources.
Darcise is working to ensure that similar principles are applied so that Indigenous people in Cameroon retain control over palm oil, which is central to their livelihoods, culture, and food systems.
In South Africa, Delwyn is building local power around food and \textbf{seeds} from the ground up.
As an activist who coordinates Greenpeace volunteers, Delwyn connects community action with broader environmental and social justice goals.
They also founded Food Connect Durban, an initiative that promotes local, sustainable food systems and the saving of heirloom \textbf{seeds}.
By encouraging people to grow, share, and protect diverse local varieties, Food Connect Durban strengthens resilience against corporate-controlled agriculture and aligns with global movements for \textbf{seed} sovereignty and food justice.
Together, these efforts form part of a wider push to reclaim food systems from corporate control and to restore power to farmers, Indigenous peoples, and local communities.
From Kenyan courts to Cameroonian classrooms and South African community gardens, activists are working to ensure that \textbf{seeds} remain a shared heritage, not a corporate asset.
Their struggles are interconnected, grounded in the belief that real food system transformation starts with defending the right to save, share, and plant \textbf{seeds} freely.

% matched lemmas: seed
\end{textsample}
